\documentclass[11pt,letterpaper]{article}
\usepackage[technical-reference]{cornell-notes}
\usepackage{needspace}

\title{Clause 11: Classes - Cornell Notes}
\author{}
\date{}
\setDocTitle{Clause 11: Classes}
\setDocSubtitle{C++ 2024 Cornell Notes}
\setDocOwner{C++ 2024 Cornell Notes}
\setDocAuthor{}
\setDocDate{}
\setCornellCollection{C++ 2024 Cornell Notes}
\setCornellUnitType{Clause}
\setCornellUnitNumber{11}
\setCornellUnitTitle{Classes}
\setCornellCompanionLabel{C++ Technical Reference}
\hypersetup{pdftitle={Clause 11: Classes - Cornell Notes},pdfsubject={C++ technical study notes},pdfauthor={}}

\begin{document}
\maketitle
\begin{CornellReferenceOrientationBox}
\textbf{Purpose.} Defines class members, special member functions, conversions, unions, inheritance, virtual dispatch, access control, construction and destruction, copy elision, and defaulted comparisons.
\par\textbf{Role.} Specifies user-defined object types and their lifetime, representation-facing, and polymorphic behavior.
\par\textbf{Principal dependencies.} Uses declarations, initialization, expressions, lookup, and object lifetime; feeds overloading and templates.
\par\textbf{Primary audience.} programmers, ABI and compiler implementers, and object-model reviewers.
\par\textbf{Major subject groups.} Preamble; Properties of classes; Class names; Class members; Unions; Local class declarations; Derived classes; Member access control; Initialization; Comparisons.
\end{CornellReferenceOrientationBox}
\section{Learning Objectives}
\begin{itemize}[label=\textcolor{CornellReferenceTeal}{$\blacktriangleright$}]
\item Define the governing ideas of Preamble and state the consequences for a program or implementation.
\item Identify the governing ideas of Properties of classes and state the consequences for a program or implementation.
\item Distinguish the governing ideas of Class names and state the consequences for a program or implementation.
\item Explain the governing ideas of Class members and state the consequences for a program or implementation.
\item Trace the governing ideas of Unions and state the consequences for a program or implementation.
\item Apply the governing ideas of Local class declarations and state the consequences for a program or implementation.
\item Analyze the governing ideas of Derived classes and state the consequences for a program or implementation.
\item Evaluate the governing ideas of Member access control and state the consequences for a program or implementation.
\item Diagnose the governing ideas of Initialization and state the consequences for a program or implementation.
\item Compare the governing ideas of Comparisons and state the consequences for a program or implementation.
\item Verify the governing ideas of General and state the consequences for a program or implementation.
\item Relate the governing ideas of Member functions and state the consequences for a program or implementation.
\item Classify the governing ideas of Non-static member functions and state the consequences for a program or implementation.
\item Justify the governing ideas of Special member functions and state the consequences for a program or implementation.
\item Audit the governing ideas of Constructors and state the consequences for a program or implementation.
\item Summarize the governing ideas of Copy/move assignment operator and state the consequences for a program or implementation.
\end{itemize}
\section{Normative Structure Map}
\small\begin{longtable}{@{}>{\RaggedRight\arraybackslash}p{0.13\textwidth}>{\RaggedRight\arraybackslash}p{0.27\textwidth}>{\RaggedRight\arraybackslash}p{0.19\textwidth}>{\RaggedRight\arraybackslash}p{0.31\textwidth}@{}}
\toprule\textbf{Subclause} & \textbf{Subject} & \textbf{Rule category} & \textbf{Key dependencies / entities}\\\midrule\endfirsthead
\toprule\textbf{Subclause} & \textbf{Subject} & \textbf{Rule category} & \textbf{Key dependencies / entities}\\\midrule\endhead
\bottomrule\endfoot
11.1 & Preamble & core-language semantics & Uses declarations, initialization, expressions, lookup, and object lifetime; feeds overloading and templates.\\
11.2 & Properties of classes & core-language semantics & Uses declarations, initialization, expressions, lookup, and object lifetime; feeds overloading and templates.\\
11.3 & Class names & core-language semantics & Uses declarations, initialization, expressions, lookup, and object lifetime; feeds overloading and templates.\\
11.4 & Class members & object contract & General, Member functions, Non-static member functions, Special member functions, and related subclauses\\
11.5 & Unions & core-language semantics & General, Anonymous unions\\
11.6 & Local class declarations & core-language semantics & Uses declarations, initialization, expressions, lookup, and object lifetime; feeds overloading and templates.\\
11.7 & Derived classes & core-language semantics & General, Multiple base classes, Virtual functions, Abstract classes\\
11.8 & Member access control & object contract & General, Access specifiers, Accessibility of base classes and base class members, Friends, and related subclauses\\
11.9 & Initialization & core-language semantics & General, Explicit initialization, Initializing bases and members, Initialization by inherited constructor, and related subclauses\\
11.10 & Comparisons & core-language semantics & Defaulted comparison operator functions, Equality operator, Three-way comparison, Secondary comparison operators\\
\end{longtable}\normalsize
\begin{CornellReferenceNormativeBox}A heading maps the clause; it does not replace the detailed conditions. For each operation, read requirements, effects, result, exceptions, complexity, remarks, and notes with their stated normative force.\end{CornellReferenceNormativeBox}
\subsection*{Rule Classification Guide}
\small\begin{longtable}{@{}>{\RaggedRight\arraybackslash}p{0.25\textwidth}>{\RaggedRight\arraybackslash}p{0.67\textwidth}@{}}\toprule\textbf{Label or category} & \textbf{How to read it}\\\midrule\endfirsthead\toprule\textbf{Label or category} & \textbf{How to read it (continued)}\\\midrule\endhead\bottomrule\endfoot
Well-formed / ill-formed & Formation is governed by syntax and semantic rules; an ill-formed program does not become well-formed because one implementation accepts it.\\
Diagnosable rule & A violation requires at least one diagnostic unless the wording places the case outside the diagnosable set.\\
No diagnostic required & The program can be ill-formed while the implementation has no diagnostic obligation.\\
Undefined behavior & No execution requirements are imposed; translation success does not establish safety or portability.\\
Unspecified behavior & One of the permitted outcomes may be chosen without documentation.\\
Implementation-defined behavior & The implementation chooses and documents the behavior.\\
Conditionally supported & The implementation may omit support but documents unsupported constructs.\\
\end{longtable}\normalsize
\section{Cornell Cue-and-Notes}
\Needspace{8\baselineskip}
\subsection*{11.1 Preamble}
\begin{longtable}{@{}>{\RaggedRight\arraybackslash\columncolor{CornellReferenceCueGray}}p{0.28\textwidth}>{\RaggedRight\arraybackslash}p{0.64\textwidth}@{}}
\rowcolor{CornellReferenceNavy}\color{white}\textbf{Cue / recall prompt} & \color{white}\textbf{Notes}\\\endfirsthead
\rowcolor{CornellReferenceNavy}\color{white}\textbf{Cue / recall prompt} & \color{white}\textbf{Notes (continued)}\\\endhead
\arrayrulecolor{CornellReferenceLineGray}
\textbf{What rule applies to Preamble?}\\[-1mm]{\scriptsize\CornellClauseRef{11.1}} & Frames the terminology, applicability, and relationships needed to interpret Preamble; detailed rules remain in the following subclauses.\\\hline
\end{longtable}
\begin{CornellReferenceSummaryBox}Preamble supplies the core-language semantics portion of this clause. The key review move is to connect its local wording to the clause-wide model, then classify each condition by normative force before relying on it.\end{CornellReferenceSummaryBox}
\Needspace{8\baselineskip}
\subsection*{11.2 Properties of classes}
\begin{longtable}{@{}>{\RaggedRight\arraybackslash\columncolor{CornellReferenceCueGray}}p{0.28\textwidth}>{\RaggedRight\arraybackslash}p{0.64\textwidth}@{}}
\rowcolor{CornellReferenceNavy}\color{white}\textbf{Cue / recall prompt} & \color{white}\textbf{Notes}\\\endfirsthead
\rowcolor{CornellReferenceNavy}\color{white}\textbf{Cue / recall prompt} & \color{white}\textbf{Notes (continued)}\\\endhead
\arrayrulecolor{CornellReferenceLineGray}
\textbf{What rule applies to Properties of classes?}\\[-1mm]{\scriptsize\CornellClauseRef{11.2}} & Defines the core-language rules for Properties of classes, including formation, semantic effect, interactions with type and lookup, and any ill-formed or implementation-dependent cases.\\\hline
\end{longtable}
\begin{CornellReferenceSummaryBox}Properties of classes supplies the core-language semantics portion of this clause. The key review move is to connect its local wording to the clause-wide model, then classify each condition by normative force before relying on it.\end{CornellReferenceSummaryBox}
\Needspace{8\baselineskip}
\subsection*{11.3 Class names}
\begin{longtable}{@{}>{\RaggedRight\arraybackslash\columncolor{CornellReferenceCueGray}}p{0.28\textwidth}>{\RaggedRight\arraybackslash}p{0.64\textwidth}@{}}
\rowcolor{CornellReferenceNavy}\color{white}\textbf{Cue / recall prompt} & \color{white}\textbf{Notes}\\\endfirsthead
\rowcolor{CornellReferenceNavy}\color{white}\textbf{Cue / recall prompt} & \color{white}\textbf{Notes (continued)}\\\endhead
\arrayrulecolor{CornellReferenceLineGray}
\textbf{What rule applies to Class names?}\\[-1mm]{\scriptsize\CornellClauseRef{11.3}} & Defines the core-language rules for Class names, including formation, semantic effect, interactions with type and lookup, and any ill-formed or implementation-dependent cases.\\\hline
\end{longtable}
\begin{CornellReferenceSummaryBox}Class names supplies the core-language semantics portion of this clause. The key review move is to connect its local wording to the clause-wide model, then classify each condition by normative force before relying on it.\end{CornellReferenceSummaryBox}
\Needspace{8\baselineskip}
\subsection*{11.4 Class members}
\begin{longtable}{@{}>{\RaggedRight\arraybackslash\columncolor{CornellReferenceCueGray}}p{0.28\textwidth}>{\RaggedRight\arraybackslash}p{0.64\textwidth}@{}}
\rowcolor{CornellReferenceNavy}\color{white}\textbf{Cue / recall prompt} & \color{white}\textbf{Notes}\\\endfirsthead
\rowcolor{CornellReferenceNavy}\color{white}\textbf{Cue / recall prompt} & \color{white}\textbf{Notes (continued)}\\\endhead
\arrayrulecolor{CornellReferenceLineGray}
\textbf{What rule applies to General?}\\[-1mm]{\scriptsize\CornellClauseRef{11.4.1}} & Frames the terminology, applicability, and relationships needed to interpret General; detailed rules remain in the following subclauses.\\\hline
\textbf{What rule applies to Member functions?}\\[-1mm]{\scriptsize\CornellClauseRef{11.4.2}} & Defines the core-language rules for Member functions, including formation, semantic effect, interactions with type and lookup, and any ill-formed or implementation-dependent cases.\\\hline
\textbf{What rule applies to Non-static member functions?}\\[-1mm]{\scriptsize\CornellClauseRef{11.4.3}} & Defines the core-language rules for Non-static member functions, including formation, semantic effect, interactions with type and lookup, and any ill-formed or implementation-dependent cases.\\\hline
\textbf{What rule applies to Special member functions?}\\[-1mm]{\scriptsize\CornellClauseRef{11.4.4}} & Defines the core-language rules for Special member functions, including formation, semantic effect, interactions with type and lookup, and any ill-formed or implementation-dependent cases.\\\hline
\textbf{How is Constructors initialized?}\\[-1mm]{\scriptsize\CornellClauseRef{11.4.5}} & Specifies how Constructors establishes object state, including participation constraints, initialization effects, and exceptional exits.\\\hline
\textbf{What rule applies to Copy/move assignment operator?}\\[-1mm]{\scriptsize\CornellClauseRef{11.4.6}} & Governs state transfer or exchange for Copy/move assignment operator; check self-operation, validity after movement, exception behavior, and invalidation.\\\hline
\textbf{What rule applies to Destructors?}\\[-1mm]{\scriptsize\CornellClauseRef{11.4.7}} & Specifies how Destructors ends the object's managed state and which cleanup or termination consequences apply.\\\hline
\textbf{When is Conversions permitted?}\\[-1mm]{\scriptsize\CornellClauseRef{11.4.8}} & Specifies source and destination conditions for Conversions, the resulting type or value category, and any information-loss, pointer, qualification, or lifetime consequence.\\\hline
\textbf{What rule applies to Static members?}\\[-1mm]{\scriptsize\CornellClauseRef{11.4.9}} & Defines the core-language rules for Static members, including formation, semantic effect, interactions with type and lookup, and any ill-formed or implementation-dependent cases.\\\hline
\textbf{What rule applies to Bit-fields?}\\[-1mm]{\scriptsize\CornellClauseRef{11.4.10}} & Defines the core-language rules for Bit-fields, including formation, semantic effect, interactions with type and lookup, and any ill-formed or implementation-dependent cases.\\\hline
\textbf{What rule applies to Allocation and deallocation functions?}\\[-1mm]{\scriptsize\CornellClauseRef{11.4.11}} & Defines the core-language rules for Allocation and deallocation functions, including formation, semantic effect, interactions with type and lookup, and any ill-formed or implementation-dependent cases.\\\hline
\textbf{What rule applies to Nested class declarations?}\\[-1mm]{\scriptsize\CornellClauseRef{11.4.12}} & Defines the core-language rules for Nested class declarations, including formation, semantic effect, interactions with type and lookup, and any ill-formed or implementation-dependent cases.\\\hline
\end{longtable}
\begin{CornellReferenceSummaryBox}Class members supplies the object contract portion of this clause. The key review move is to connect its local wording to the clause-wide model, then classify each condition by normative force before relying on it.\end{CornellReferenceSummaryBox}
\Needspace{5\baselineskip}\paragraph{Detailed coverage ledger.}
\begin{description}[style=nextline,leftmargin=2.8cm,font=\normalfont\bfseries\color{CornellReferenceTeal}]
\item[11.4.5.1] \textbf{General.} Frames the terminology, applicability, and relationships needed to interpret General; detailed rules remain in the following subclauses.
\item[11.4.5.2] \textbf{Default constructors.} Specifies how Default constructors establishes object state, including participation constraints, initialization effects, and exceptional exits.
\item[11.4.5.3] \textbf{Copy/move constructors.} Specifies how Copy/move constructors establishes object state, including participation constraints, initialization effects, and exceptional exits.
\item[11.4.8.1] \textbf{General.} Frames the terminology, applicability, and relationships needed to interpret General; detailed rules remain in the following subclauses.
\item[11.4.8.2] \textbf{Conversion by constructor.} Specifies how Conversion by constructor establishes object state, including participation constraints, initialization effects, and exceptional exits.
\item[11.4.8.3] \textbf{Conversion functions.} Specifies source and destination conditions for Conversion functions, the resulting type or value category, and any information-loss, pointer, qualification, or lifetime consequence.
\item[11.4.9.1] \textbf{General.} Frames the terminology, applicability, and relationships needed to interpret General; detailed rules remain in the following subclauses.
\item[11.4.9.2] \textbf{Static member functions.} Defines the core-language rules for Static member functions, including formation, semantic effect, interactions with type and lookup, and any ill-formed or implementation-dependent cases.
\item[11.4.9.3] \textbf{Static data members.} Defines the core-language rules for Static data members, including formation, semantic effect, interactions with type and lookup, and any ill-formed or implementation-dependent cases.
\end{description}
\Needspace{8\baselineskip}
\subsection*{11.5 Unions}
\begin{longtable}{@{}>{\RaggedRight\arraybackslash\columncolor{CornellReferenceCueGray}}p{0.28\textwidth}>{\RaggedRight\arraybackslash}p{0.64\textwidth}@{}}
\rowcolor{CornellReferenceNavy}\color{white}\textbf{Cue / recall prompt} & \color{white}\textbf{Notes}\\\endfirsthead
\rowcolor{CornellReferenceNavy}\color{white}\textbf{Cue / recall prompt} & \color{white}\textbf{Notes (continued)}\\\endhead
\arrayrulecolor{CornellReferenceLineGray}
\textbf{What rule applies to General?}\\[-1mm]{\scriptsize\CornellClauseRef{11.5.1}} & Frames the terminology, applicability, and relationships needed to interpret General; detailed rules remain in the following subclauses.\\\hline
\textbf{What rule applies to Anonymous unions?}\\[-1mm]{\scriptsize\CornellClauseRef{11.5.2}} & Defines the core-language rules for Anonymous unions, including formation, semantic effect, interactions with type and lookup, and any ill-formed or implementation-dependent cases.\\\hline
\end{longtable}
\begin{CornellReferenceSummaryBox}Unions supplies the core-language semantics portion of this clause. The key review move is to connect its local wording to the clause-wide model, then classify each condition by normative force before relying on it.\end{CornellReferenceSummaryBox}
\Needspace{8\baselineskip}
\subsection*{11.6 Local class declarations}
\begin{longtable}{@{}>{\RaggedRight\arraybackslash\columncolor{CornellReferenceCueGray}}p{0.28\textwidth}>{\RaggedRight\arraybackslash}p{0.64\textwidth}@{}}
\rowcolor{CornellReferenceNavy}\color{white}\textbf{Cue / recall prompt} & \color{white}\textbf{Notes}\\\endfirsthead
\rowcolor{CornellReferenceNavy}\color{white}\textbf{Cue / recall prompt} & \color{white}\textbf{Notes (continued)}\\\endhead
\arrayrulecolor{CornellReferenceLineGray}
\textbf{What rule applies to Local class declarations?}\\[-1mm]{\scriptsize\CornellClauseRef{11.6}} & Defines the core-language rules for Local class declarations, including formation, semantic effect, interactions with type and lookup, and any ill-formed or implementation-dependent cases.\\\hline
\end{longtable}
\begin{CornellReferenceSummaryBox}Local class declarations supplies the core-language semantics portion of this clause. The key review move is to connect its local wording to the clause-wide model, then classify each condition by normative force before relying on it.\end{CornellReferenceSummaryBox}
\Needspace{8\baselineskip}
\subsection*{11.7 Derived classes}
\begin{longtable}{@{}>{\RaggedRight\arraybackslash\columncolor{CornellReferenceCueGray}}p{0.28\textwidth}>{\RaggedRight\arraybackslash}p{0.64\textwidth}@{}}
\rowcolor{CornellReferenceNavy}\color{white}\textbf{Cue / recall prompt} & \color{white}\textbf{Notes}\\\endfirsthead
\rowcolor{CornellReferenceNavy}\color{white}\textbf{Cue / recall prompt} & \color{white}\textbf{Notes (continued)}\\\endhead
\arrayrulecolor{CornellReferenceLineGray}
\textbf{What rule applies to General?}\\[-1mm]{\scriptsize\CornellClauseRef{11.7.1}} & Frames the terminology, applicability, and relationships needed to interpret General; detailed rules remain in the following subclauses.\\\hline
\textbf{What rule applies to Multiple base classes?}\\[-1mm]{\scriptsize\CornellClauseRef{11.7.2}} & Defines the core-language rules for Multiple base classes, including formation, semantic effect, interactions with type and lookup, and any ill-formed or implementation-dependent cases.\\\hline
\textbf{What rule applies to Virtual functions?}\\[-1mm]{\scriptsize\CornellClauseRef{11.7.3}} & Defines the core-language rules for Virtual functions, including formation, semantic effect, interactions with type and lookup, and any ill-formed or implementation-dependent cases.\\\hline
\textbf{What rule applies to Abstract classes?}\\[-1mm]{\scriptsize\CornellClauseRef{11.7.4}} & Defines the core-language rules for Abstract classes, including formation, semantic effect, interactions with type and lookup, and any ill-formed or implementation-dependent cases.\\\hline
\end{longtable}
\begin{CornellReferenceSummaryBox}Derived classes supplies the core-language semantics portion of this clause. The key review move is to connect its local wording to the clause-wide model, then classify each condition by normative force before relying on it.\end{CornellReferenceSummaryBox}
\Needspace{8\baselineskip}
\subsection*{11.8 Member access control}
\begin{longtable}{@{}>{\RaggedRight\arraybackslash\columncolor{CornellReferenceCueGray}}p{0.28\textwidth}>{\RaggedRight\arraybackslash}p{0.64\textwidth}@{}}
\rowcolor{CornellReferenceNavy}\color{white}\textbf{Cue / recall prompt} & \color{white}\textbf{Notes}\\\endfirsthead
\rowcolor{CornellReferenceNavy}\color{white}\textbf{Cue / recall prompt} & \color{white}\textbf{Notes (continued)}\\\endhead
\arrayrulecolor{CornellReferenceLineGray}
\textbf{What rule applies to General?}\\[-1mm]{\scriptsize\CornellClauseRef{11.8.1}} & Frames the terminology, applicability, and relationships needed to interpret General; detailed rules remain in the following subclauses.\\\hline
\textbf{What rule applies to Access specifiers?}\\[-1mm]{\scriptsize\CornellClauseRef{11.8.2}} & Defines the core-language rules for Access specifiers, including formation, semantic effect, interactions with type and lookup, and any ill-formed or implementation-dependent cases.\\\hline
\textbf{What rule applies to Accessibility of base classes and base class members?}\\[-1mm]{\scriptsize\CornellClauseRef{11.8.3}} & Defines the core-language rules for Accessibility of base classes and base class members, including formation, semantic effect, interactions with type and lookup, and any ill-formed or implementation-dependent cases.\\\hline
\textbf{What rule applies to Friends?}\\[-1mm]{\scriptsize\CornellClauseRef{11.8.4}} & Defines the core-language rules for Friends, including formation, semantic effect, interactions with type and lookup, and any ill-formed or implementation-dependent cases.\\\hline
\textbf{What rule applies to Protected member access?}\\[-1mm]{\scriptsize\CornellClauseRef{11.8.5}} & Defines the core-language rules for Protected member access, including formation, semantic effect, interactions with type and lookup, and any ill-formed or implementation-dependent cases.\\\hline
\textbf{What rule applies to Access to virtual functions?}\\[-1mm]{\scriptsize\CornellClauseRef{11.8.6}} & Defines the core-language rules for Access to virtual functions, including formation, semantic effect, interactions with type and lookup, and any ill-formed or implementation-dependent cases.\\\hline
\textbf{What rule applies to Multiple access?}\\[-1mm]{\scriptsize\CornellClauseRef{11.8.7}} & Defines the core-language rules for Multiple access, including formation, semantic effect, interactions with type and lookup, and any ill-formed or implementation-dependent cases.\\\hline
\textbf{What rule applies to Nested classes?}\\[-1mm]{\scriptsize\CornellClauseRef{11.8.8}} & Defines the core-language rules for Nested classes, including formation, semantic effect, interactions with type and lookup, and any ill-formed or implementation-dependent cases.\\\hline
\end{longtable}
\begin{CornellReferenceSummaryBox}Member access control supplies the object contract portion of this clause. The key review move is to connect its local wording to the clause-wide model, then classify each condition by normative force before relying on it.\end{CornellReferenceSummaryBox}
\Needspace{8\baselineskip}
\subsection*{11.9 Initialization}
\begin{longtable}{@{}>{\RaggedRight\arraybackslash\columncolor{CornellReferenceCueGray}}p{0.28\textwidth}>{\RaggedRight\arraybackslash}p{0.64\textwidth}@{}}
\rowcolor{CornellReferenceNavy}\color{white}\textbf{Cue / recall prompt} & \color{white}\textbf{Notes}\\\endfirsthead
\rowcolor{CornellReferenceNavy}\color{white}\textbf{Cue / recall prompt} & \color{white}\textbf{Notes (continued)}\\\endhead
\arrayrulecolor{CornellReferenceLineGray}
\textbf{What rule applies to General?}\\[-1mm]{\scriptsize\CornellClauseRef{11.9.1}} & Frames the terminology, applicability, and relationships needed to interpret General; detailed rules remain in the following subclauses.\\\hline
\textbf{What rule applies to Explicit initialization?}\\[-1mm]{\scriptsize\CornellClauseRef{11.9.2}} & Determines the initialization form used by Explicit initialization, the candidates or conversions considered, object-lifetime start, narrowing rules, and failure consequences.\\\hline
\textbf{What rule applies to Initializing bases and members?}\\[-1mm]{\scriptsize\CornellClauseRef{11.9.3}} & Determines the initialization form used by Initializing bases and members, the candidates or conversions considered, object-lifetime start, narrowing rules, and failure consequences.\\\hline
\textbf{How is Initialization by inherited constructor initialized?}\\[-1mm]{\scriptsize\CornellClauseRef{11.9.4}} & Specifies how Initialization by inherited constructor establishes object state, including participation constraints, initialization effects, and exceptional exits.\\\hline
\textbf{What rule applies to Construction and destruction?}\\[-1mm]{\scriptsize\CornellClauseRef{11.9.5}} & Defines the core-language rules for Construction and destruction, including formation, semantic effect, interactions with type and lookup, and any ill-formed or implementation-dependent cases.\\\hline
\textbf{What rule applies to Copy/move elision?}\\[-1mm]{\scriptsize\CornellClauseRef{11.9.6}} & Defines the core-language rules for Copy/move elision, including formation, semantic effect, interactions with type and lookup, and any ill-formed or implementation-dependent cases.\\\hline
\end{longtable}
\begin{CornellReferenceSummaryBox}Initialization supplies the core-language semantics portion of this clause. The key review move is to connect its local wording to the clause-wide model, then classify each condition by normative force before relying on it.\end{CornellReferenceSummaryBox}
\Needspace{8\baselineskip}
\subsection*{11.10 Comparisons}
\begin{longtable}{@{}>{\RaggedRight\arraybackslash\columncolor{CornellReferenceCueGray}}p{0.28\textwidth}>{\RaggedRight\arraybackslash}p{0.64\textwidth}@{}}
\rowcolor{CornellReferenceNavy}\color{white}\textbf{Cue / recall prompt} & \color{white}\textbf{Notes}\\\endfirsthead
\rowcolor{CornellReferenceNavy}\color{white}\textbf{Cue / recall prompt} & \color{white}\textbf{Notes (continued)}\\\endhead
\arrayrulecolor{CornellReferenceLineGray}
\textbf{What rule applies to Defaulted comparison operator functions?}\\[-1mm]{\scriptsize\CornellClauseRef{11.10.1}} & Defines the core-language rules for Defaulted comparison operator functions, including formation, semantic effect, interactions with type and lookup, and any ill-formed or implementation-dependent cases.\\\hline
\textbf{What rule applies to Equality operator?}\\[-1mm]{\scriptsize\CornellClauseRef{11.10.2}} & Defines the core-language rules for Equality operator, including formation, semantic effect, interactions with type and lookup, and any ill-formed or implementation-dependent cases.\\\hline
\textbf{What rule applies to Three-way comparison?}\\[-1mm]{\scriptsize\CornellClauseRef{11.10.3}} & Defines the core-language rules for Three-way comparison, including formation, semantic effect, interactions with type and lookup, and any ill-formed or implementation-dependent cases.\\\hline
\textbf{What rule applies to Secondary comparison operators?}\\[-1mm]{\scriptsize\CornellClauseRef{11.10.4}} & Defines the core-language rules for Secondary comparison operators, including formation, semantic effect, interactions with type and lookup, and any ill-formed or implementation-dependent cases.\\\hline
\end{longtable}
\begin{CornellReferenceSummaryBox}Comparisons supplies the core-language semantics portion of this clause. The key review move is to connect its local wording to the clause-wide model, then classify each condition by normative force before relying on it.\end{CornellReferenceSummaryBox}
\section{Language-Lawyer and Library-Usage Pitfalls}
\begin{CornellReferencePitfallBox}\begin{itemize}[label=\textcolor{CornellReferenceAmber}{$\blacktriangleright$}]
\item Assuming an implicitly declared special member is necessarily defined or usable.
\item Calling virtual behavior during construction as though the most-derived object were active.
\item Confusing access control with name lookup.
\item Ignoring active-member and lifetime rules in unions.
\item Expecting copy elision in a context where it is only permitted, not guaranteed.
\item Treating a note, example, or recommended practice as though it imposed a mandatory program rule.
\end{itemize}\end{CornellReferencePitfallBox}
\section{Programmer and Implementation Checklist}
\begin{CornellReferenceChecklistBox}\begin{itemize}[label=$\square$]
\item Determine which special members are declared, deleted, trivial, or eligible.
\item Track base and member initialization order.
\item Find the final overrider for each virtual call context.
\item Audit access after lookup has identified a declaration.
\item Check union active-member transitions and lifetime.
\item Verify conversion functions and converting constructors.
\item Evaluate generated comparison category and member order.
\item For \CornellClauseRef{11.1}, verify preamble against the precise rule category and all cited dependencies.
\item For \CornellClauseRef{11.2}, verify properties of classes against the precise rule category and all cited dependencies.
\item For \CornellClauseRef{11.3}, verify class names against the precise rule category and all cited dependencies.
\item For \CornellClauseRef{11.4}, verify class members against the precise rule category and all cited dependencies.
\end{itemize}\end{CornellReferenceChecklistBox}
\section{Clause Synthesis}
\begin{CornellReferenceOrientationBox}Defines class members, special member functions, conversions, unions, inheritance, virtual dispatch, access control, construction and destruction, copy elision, and defaulted comparisons. Specifies user-defined object types and their lifetime, representation-facing, and polymorphic behavior. A sound review distinguishes syntax and participation from runtime preconditions, keeps notes and examples non-mandatory, and follows cross-clause dependencies before making a portability claim.\end{CornellReferenceOrientationBox}
\section{Self-Test Questions}
\begin{enumerate}
\item What role does Preamble play in the clause's overall model?
\item Which normative distinctions must be kept separate when applying Properties of classes?
\item How would you audit a program or implementation that depends on Class names?
\item Compare Class members with Unions; what different question does each answer?
\item What role does Unions play in the clause's overall model?
\item Which normative distinctions must be kept separate when applying Local class declarations?
\item How would you audit a program or implementation that depends on Derived classes?
\item Compare Member access control with Initialization; what different question does each answer?
\item What role does Initialization play in the clause's overall model?
\item Which normative distinctions must be kept separate when applying Comparisons?
\item How would you audit a program or implementation that depends on General?
\item Compare Member functions with Non-static member functions; what different question does each answer?
\item What role does Non-static member functions play in the clause's overall model?
\item Which normative distinctions must be kept separate when applying Special member functions?
\item Scenario: a program relies on an unstated implementation behavior while using Constructors. What must be checked before calling the program portable?
\item Scenario: a program relies on an unstated implementation behavior while using Copy/move assignment operator. What must be checked before calling the program portable?
\end{enumerate}
\clearpage\section{Answer Key}
\begin{enumerate}
\item Preamble is one part of the clause's core-language semantics model. Its role is understood by connecting the local rule in 11.1 to the clause orientation and its dependent concepts.
\item Keep formation and well-formedness separate from runtime preconditions and effects. Also distinguish mandatory wording from notes, implementation choice from undefined behavior, and interface declarations from detailed contracts; see 11.2.
\item Audit Class names by locating the controlling wording in 11.3, listing inputs and type constraints, proving any preconditions, recording effects and failure behavior, and checking lifetime, invalidation, complexity, and synchronization where applicable.
\item Class members and Unions occupy different steps or facility groups. Analyze each under its own subclause, then follow the explicit dependency between them rather than transferring one set of guarantees to the other; begin at 11.4.
\item Unions is one part of the clause's core-language semantics model. Its role is understood by connecting the local rule in 11.5 to the clause orientation and its dependent concepts.
\item Keep formation and well-formedness separate from runtime preconditions and effects. Also distinguish mandatory wording from notes, implementation choice from undefined behavior, and interface declarations from detailed contracts; see 11.6.
\item Audit Derived classes by locating the controlling wording in 11.7, listing inputs and type constraints, proving any preconditions, recording effects and failure behavior, and checking lifetime, invalidation, complexity, and synchronization where applicable.
\item Member access control and Initialization occupy different steps or facility groups. Analyze each under its own subclause, then follow the explicit dependency between them rather than transferring one set of guarantees to the other; begin at 11.8.
\item Initialization is one part of the clause's core-language semantics model. Its role is understood by connecting the local rule in 11.9 to the clause orientation and its dependent concepts.
\item Keep formation and well-formedness separate from runtime preconditions and effects. Also distinguish mandatory wording from notes, implementation choice from undefined behavior, and interface declarations from detailed contracts; see 11.10.
\item Audit General by locating the controlling wording in 11.9.1, listing inputs and type constraints, proving any preconditions, recording effects and failure behavior, and checking lifetime, invalidation, complexity, and synchronization where applicable.
\item Member functions and Non-static member functions occupy different steps or facility groups. Analyze each under its own subclause, then follow the explicit dependency between them rather than transferring one set of guarantees to the other; begin at 11.4.2.
\item Non-static member functions is one part of the clause's object contract model. Its role is understood by connecting the local rule in 11.4.3 to the clause orientation and its dependent concepts.
\item Keep formation and well-formedness separate from runtime preconditions and effects. Also distinguish mandatory wording from notes, implementation choice from undefined behavior, and interface declarations from detailed contracts; see 11.4.4.
\item First identify the exact requirement in 11.4.5, then classify the relied-on behavior as required, unspecified, implementation-defined, conditionally supported, or undefined. Portability requires satisfying the program obligations and avoiding dependence on a choice not guaranteed in the target environments.
\item First identify the exact requirement in 11.4.6, then classify the relied-on behavior as required, unspecified, implementation-defined, conditionally supported, or undefined. Portability requires satisfying the program obligations and avoiding dependence on a choice not guaranteed in the target environments.
\end{enumerate}
\section{Key-Term Recap}
\begin{longtable}{@{}>{\RaggedRight\arraybackslash}p{0.28\textwidth}>{\RaggedRight\arraybackslash}p{0.64\textwidth}@{}}\toprule\textbf{Term} & \textbf{Working meaning}\\\midrule\endfirsthead\toprule\textbf{Term} & \textbf{Working meaning (continued)}\\\midrule\endhead\bottomrule\endfoot
special member function & A constructor, assignment operator, or destructor subject to implicit declaration and generation rules.\\
virtual function & A member function dispatched through the final overrider for an object.\\
base class & A class incorporated through inheritance into a derived class.\\
abstract class & A class that cannot be directly instantiated because required virtual behavior remains pure.\\
union & A class whose non-static data members share storage under active-member rules.\\
copy elision & Specified omission of copying or moving in designated initialization contexts.\\
defaulted comparison & A comparison whose generated behavior follows the member and base subobject model.\\
\end{longtable}
\CornellTakeaway{Class correctness depends on the exact interaction of member generation, initialization order, lifetime, inheritance, access, and dispatch.}
\end{document}
